\documentclass[10pt,a4paper]{amsart}
\usepackage{amsmath,amssymb,amsthm,graphicx}
% Retained minimal preamble for the Stacks Project chapter more-morphisms.tex (More on Morphisms).
% Verbatim excerpt of the chapter preamble: Xy-pic setup (native preamble lines 14-18) and the native
% theorem declarations with the mathematical macros (native preamble lines 154-217). Copied mechanically;
% nothing added, removed or redefined. The wrapper class and the amsmath/amssymb/amsthm/graphicx base
% packages are supplied by the assembler, so the class/externaldocument part of the native preamble is omitted.
% For commutative diagrams we use Xy-pic
\usepackage[all]{xy}

% We use 2cell for 2-commutative diagrams.
\xyoption{2cell}
% Theorem environments.
%
\theoremstyle{plain}
\newtheorem{theorem}[subsection]{Theorem}
\newtheorem{proposition}[subsection]{Proposition}
\newtheorem{lemma}[subsection]{Lemma}

\theoremstyle{definition}
\newtheorem{definition}[subsection]{Definition}
\newtheorem{example}[subsection]{Example}
\newtheorem{exercise}[subsection]{Exercise}
\newtheorem{situation}[subsection]{Situation}

\theoremstyle{remark}
\newtheorem{remark}[subsection]{Remark}
\newtheorem{remarks}[subsection]{Remarks}

\numberwithin{equation}{subsection}

% Macros
%
\def\lim{\mathop{\mathrm{lim}}\nolimits}
\def\colim{\mathop{\mathrm{colim}}\nolimits}
\def\Spec{\mathop{\mathrm{Spec}}}
\def\Hom{\mathop{\mathrm{Hom}}\nolimits}
\def\Ext{\mathop{\mathrm{Ext}}\nolimits}
\def\SheafHom{\mathop{\mathcal{H}\!\mathit{om}}\nolimits}
\def\SheafExt{\mathop{\mathcal{E}\!\mathit{xt}}\nolimits}
\def\Sch{\mathit{Sch}}
\def\Mor{\mathop{\mathrm{Mor}}\nolimits}
\def\Ob{\mathop{\mathrm{Ob}}\nolimits}
\def\Sh{\mathop{\mathit{Sh}}\nolimits}
\def\NL{\mathop{N\!L}\nolimits}
\def\CH{\mathop{\mathrm{CH}}\nolimits}
\def\proetale{{pro\text{-}\acute{e}tale}}
\def\etale{{\acute{e}tale}}
\def\QCoh{\mathit{QCoh}}
\def\Ker{\mathop{\mathrm{Ker}}}
\def\Im{\mathop{\mathrm{Im}}}
\def\Coker{\mathop{\mathrm{Coker}}}
\def\Coim{\mathop{\mathrm{Coim}}}

% Boxtimes
%
\DeclareMathSymbol{\boxtimes}{\mathbin}{AMSa}{"02}

%
% Macros for moduli stacks/spaces
%
\def\QCohstack{\mathcal{QC}\!\mathit{oh}}
\def\Cohstack{\mathcal{C}\!\mathit{oh}}
\def\Spacesstack{\mathcal{S}\!\mathit{paces}}
\def\Quotfunctor{\mathrm{Quot}}
\def\Hilbfunctor{\mathrm{Hilb}}
\def\Curvesstack{\mathcal{C}\!\mathit{urves}}
\def\Polarizedstack{\mathcal{P}\!\mathit{olarized}}
\def\Complexesstack{\mathcal{C}\!\mathit{omplexes}}
% \Pic is the operator that assigns to X its picard group, usage \Pic(X)
% \Picardstack_{X/B} denotes the Picard stack of X over B
% \Picardfunctor_{X/B} denotes the Picard functor of X over B
\def\Pic{\mathop{\mathrm{Pic}}\nolimits}
\def\Picardstack{\mathcal{P}\!\mathit{ic}}
\def\Picardfunctor{\mathrm{Pic}}
\def\Deformationcategory{\mathcal{D}\!\mathit{ef}}
\usepackage{cleveref}
\crefname{originaltheorem}{original theorem}{original theorems}
\Crefname{originaltheorem}{Original theorem}{Original theorems}
\makeatletter
\newlabel{theorem-normalized-base-change-with-reduced-fibre}{{09IL}{source}{}{sourcepointer-1}{}}
\newlabel{theorem-normalized-base-change-with-reduced-fibre@cref}{{[originaltheorem][0][]09IL}{[1][0][]source}{}{}{}}
\makeatother
\title{Normalized base change over a one-dimensional regular base is smooth at generic fibre points}
\author{The Stacks Project Authors}
\date{stacks-project:more-morphisms.tex:local-snapshot-20261011; GNU FDL 1.2 or later}
\begin{document}
\maketitle
\noindent\textit{Editorial note.} Counted span is the author's own proof environment, source lines 19374-19398, immediately following the statement at 19357-19372.\par\medskip
\noindent\textit{Editorial note.} The proof invokes the general theorem twice by \textbackslash{}ref; those are navigation pointers into the same chapter, and the smoothness locus transported here is the one displayed in the span.\par\medskip
\noindent\textit{Editorial note.} The author's phrase 'in the proof of Theorem ... we have seen that' refers to a computation inside the cited theorem; it is retained verbatim.\par\medskip
\noindent\textit{Allowed standard prerequisites (editorial).}\begin{itemize}
\item Normalized base change with reduced fibre, general form. chapter theorem theorem-normalized-base-change-with-reduced-fibre, cited at proof lines 19379 and 19394. Supplies, over each affine chart of S, the finite extension L\_i/K and the smoothness statement whose locus is transported to the global situation.
\item Normalization of a scheme in a finite field extension. standard Stacks construction used at proof lines 19369-19371 and 19381-19382. Defines T and the normalized base change Y over the integral closure B\_i of A\_i in L\_i.
\item Finite extensions of a field are directed and admit a common dominating extension. used at proof line 19383. Produces a single finite L/K dominating all the chartwise extensions L\_i/K.
\end{itemize}\medskip
\ifdefined\hypertarget\hypertarget{sourcepointer-1}{}\fi
\noindent\textit{Source reference.} Original theorem 09IL: sources/stacks-more-morphisms/pinned-original-source.tex, source lines 19282--19282. Navigation pointer only: the counted proof cites this statement by \textbackslash{}ref as a standing input; the deduction the counted proof has to make is carried out inside the counted span, and no part of that deduction is delegated to this pointer.\par\medskip
\begin{lemma}[Variant over curves]
\label{lemma-normalized-base-change-with-reduced-fibre-over-curve}
Let $f : X \to S$ be a flat, finite type morphism of schemes.
Assume $S$ is Nagata, integral with function field $K$, and
regular of dimension $1$. Then there exists a finite extension $L/K$
such that in the diagram
$$
\xymatrix{
Y \ar[rd]_g \ar[r]_-\nu & X \times_S T \ar[d] \ar[r] & X \ar[d]_f \\
& T \ar[r] & S
}
$$
the morphism $g$ is smooth at all generic points of fibres. Here
$T$ is the normalization of $S$ in $\Spec(L)$ and $\nu : Y \to X \times_S T$
is the normalization.
\end{lemma}

\begin{proof}
Choose a finite affine open covering $S = \bigcup \Spec(A_i)$.
Then $K$ is equal to the fraction field of $A_i$ for all $i$.
Let $X_i = X \times_S \Spec(A_i)$.
Choose $L_i/K$ as in
Theorem \ref{theorem-normalized-base-change-with-reduced-fibre}
for the morphism $X_i \to \Spec(A_i)$.
Let $B_i \subset L_i$ be the integral closure of $A_i$ and
let $Y_i$ be the normalized base change of $X$ to $B_i$.
Let $L/K$ be a finite extension dominating each $L_i$.
Let $T_i \subset T$ be the inverse image of $\Spec(A_i)$.
For each $i$ we get a commutative diagram
$$
\xymatrix{
g^{-1}(T_i) \ar[r] \ar[d] & Y_i \ar[r] \ar[d] & X \times_S \Spec(A_i) \ar[d] \\
T_i \ar[r] & \Spec(B_i) \ar[r] & \Spec(A_i)
}
$$
and in fact the left hand square is a normalized base change
as discussed at the beginning of the section. In the proof
of Theorem \ref{theorem-normalized-base-change-with-reduced-fibre}
we have seen that the smooth locus of $Y \to T$ contains the
inverse image in $g^{-1}(T_i)$ of the set of points
where $Y_i$ is smooth over $B_i$. This proves the lemma.
\end{proof}

\end{document}
